
\chapter{主存、总线、I/O 与 DMA}

前面已经解释了最小 CPU 如何在内部取指、译码、执行和提交状态。若停在这里，读者只能得到一台“会在寄存器之间搬数”的小机器。真正的整机还需要大容量主存、可寻址设备、外部事件、等待、错误和批量数据搬运。本章专门补上这条桥：CPU 怎样把一个地址变成一次存储器或设备事务，设备怎样把完成事件通知 CPU，DMA 又怎样在不让 CPU 逐字搬运的情况下读写主存。

本章的目标不是背总线协议名字。无论是 8086 板级总线、微控制器片内总线、PCIe 设备，还是现代 SoC 互连，工程问题都相同：谁发起请求，谁响应地址，数据由谁驱动，什么时候采样，目标慢了怎样等待，目标不存在怎样失败，设备完成后怎样通知。把这些问题写清楚，才能从最小 CPU 走到可启动、可调试、可扩展的整机。

\begin{keyidea}
整机硬件不是“CPU 加一块内存”的粗线图，而是一组可验收事务。每次访问都必须说明地址、方向、字节选择、目标片选、握手、完成、错误和副作用。只有这些边界明确，CPU 才能可靠地取指、访存、访问设备和处理中断。
\end{keyidea}

\begin{pdfdiagram}
  \node[hw state, minimum width=1.55cm] (cpu) at (0,0) {CPU};
  \node[hw compute, minimum width=1.95cm] (dec) at (2.35,0) {地址译码};
  \node[hw state, minimum width=1.65cm] (rom) at (5.0,1.0) {ROM};
  \node[hw state, minimum width=1.65cm] (ram) at (5.0,0) {SRAM};
  \node[hw control, minimum width=1.65cm] (io) at (5.0,-1.0) {MMIO};
  \node[hw control, minimum width=1.75cm] (irq) at (7.45,-1.0) {中断/DMA};
  \draw[hw bus] (cpu) -- node[above,hw label]{地址/控制} (dec);
  \draw[hw bus] (dec.east) |- (rom.west) node[pos=0.25,above,hw label]{片选};
  \draw[hw bus] (dec) -- (ram);
  \draw[hw bus] (dec.east) |- (io.west);
  \draw[hw return] (io) -- (irq);
  \draw[hw return] (irq.north) |- (cpu.south) node[pos=0.2,below,hw label]{完成/事件};
  \node[hw label, text width=8.8cm] at (3.75,-2.0) {整机事务从 CPU 发起，但完成、等待和错误由目标与互连共同决定；设备还可通过中断或 DMA 把事件带回 CPU。};
\end{pdfdiagram}
\diagramnote{从最小 CPU 到整机事务。地址译码决定目标，握手和事件决定本次访问何时真正完成。}

\section{一、主存不是抽象数组}

软件把内存看成按地址编号的字节序列，硬件却必须用地址线、译码器、存储阵列、数据线和控制信号实现这件事。一次读访问至少包含四个事实：发起者给出地址，译码逻辑选中目标，目标经过内部访问后驱动数据，发起者在规定边界采样。一次写访问则要求地址、写数据、字节使能和写控制在目标采样窗口内稳定。

\begin{longtable}{L{0.18\linewidth}L{0.34\linewidth}L{0.38\linewidth}}
\toprule
对象 & 硬件含义 & 验收问题 \\
\midrule
地址 & 指出本次事务目标位置 & 高位译码是否只选中一个目标 \\
数据 & 读时由目标返回，写时由发起者提供 & 同一时刻是否只有一个驱动者 \\
读写控制 & 区分本次事务方向和类型 & 输出使能和写使能是否互斥 \\
字节使能 & 指明哪些 8-bit 车道有效 & 字节写不会破坏相邻数据 \\
ready/valid & 指明请求或响应是否真正完成 & 慢目标能否插入等待而不丢地址和数据 \\
\bottomrule
\end{longtable}

\topic{ROM、SRAM 和 DRAM 的读写不是同一种动作}
ROM 或 Flash 适合保存启动代码，掉电后内容仍在；SRAM 接口直观，地址稳定后在访问时间内返回数据，适合教学板、cache 和小容量片上 RAM；DRAM 密度高，但单元需要刷新，访问还要经过 bank、行、列、预充电和控制器调度。把这些对象都叫“内存”没有错，但在硬件验收时必须区分它们的时序和控制边界。

\begin{longtable}{L{0.18\linewidth}L{0.34\linewidth}L{0.38\linewidth}}
\toprule
存储类型 & 典型用途 & 关键边界 \\
\midrule
ROM/Flash & 复位入口、固件、常量表 & 复位地址必须映射到有效内容，读时序满足取指 \\
SRAM & cache、实验板 RAM、小容量片上存储 & \code{CS}/\code{OE}/\code{WE} 互斥，地址和数据满足建立时间 \\
DRAM & 大容量主存 & 控制器负责 activate、read/write、precharge、refresh \\
寄存器文件 & CPU 内部近距离状态 & 多读端口、写端口、旁路和同周期读写规则明确 \\
\bottomrule
\end{longtable}

\topic{异步 SRAM 是最好的第一块主存}
教学整机可以先使用低速异步 SRAM。读周期中，地址和片选稳定后，打开输出使能，SRAM 在访问时间后驱动数据线；写周期中，输出使能关闭，发起者驱动数据线，在写使能有效窗口内把数据写入目标字节。这个模型足够简单，也足以暴露硬件最常见错误：片选重叠、输出争用、写使能过早、字节使能错误。

\begin{lstlisting}
SRAM read:
  address stable
  CS active, OE active, WE inactive
  wait access time
  sample data

SRAM write:
  address and write_data stable
  CS active, OE inactive, WE active for write window
  selected byte lanes store data
\end{lstlisting}

\begin{pdfdiagram}[x=1cm,y=1cm]
  \draw[BookRule,line width=0.45pt] (0,1.4) -- (9.6,1.4);
  \draw[BookRule,line width=0.45pt] (0,0.65) -- (9.6,0.65);
  \draw[BookRule,line width=0.45pt] (0,-0.1) -- (9.6,-0.1);
  \node[left,hw label] at (0,1.4) {ADDR};
  \node[left,hw label] at (0,0.65) {CS/OE};
  \node[left,hw label] at (0,-0.1) {DATA};
  \draw[BookTeal,line width=0.85pt] (0.8,1.4) -- (8.8,1.4);
  \draw[BookAccent,line width=0.85pt] (1.4,0.65) -- (7.9,0.65);
  \draw[BookMuted,line width=0.65pt,dashed] (0.8,-0.1) -- (4.35,-0.1);
  \draw[BookTeal,line width=0.85pt] (4.35,-0.1) -- (8.2,-0.1);
  \draw[BookRule,line width=0.45pt] (0.8,1.75) -- (0.8,-0.45);
  \draw[BookRule,line width=0.45pt] (4.35,1.75) -- (4.35,-0.45);
  \draw[BookRule,line width=0.45pt] (8.2,1.75) -- (8.2,-0.45);
  \node[hw label] at (0.8,1.98) {地址稳定};
  \node[hw label] at (4.35,1.98) {数据有效};
  \node[hw label] at (8.2,1.98) {CPU 采样};
  \node[hw label, text width=8.6cm] at (4.8,-0.9) {等待状态的目的，是把采样点推迟到数据有效之后，同时保持地址、片选和方向不变。};
\end{pdfdiagram}
\diagramnote{异步 SRAM 读周期的最小账本。数据有效之前不能被 CPU 采样。}

\topic{DRAM 需要控制器而不是直接接地址线}
DRAM 的容量来自更密集的存储单元，代价是接口复杂。控制器要先打开某个 bank 的某一行，再读写列；换行前需要预充电；即使没有普通访问，也要周期性刷新；高速 DRAM 还需要训练采样窗口。CPU 看到的是“主存读写”，但实际路径经过控制器、队列、bank 调度、行命中或行冲突。这也是为什么现代 CPU 必须靠 cache 把常用数据放近。

\begin{longtable}{L{0.2\linewidth}L{0.34\linewidth}L{0.36\linewidth}}
\toprule
DRAM 动作 & 硬件含义 & 时延来源 \\
\midrule
Activate & 打开某个 bank 中的一行 & 行译码和感放建立 \\
Read/Write & 在已打开行内选择列传输数据 & 列访问、突发传输和总线排队 \\
Precharge & 关闭当前行，准备访问另一行 & 位线恢复到可再次访问状态 \\
Refresh & 周期性恢复存储单元电荷 & 控制器必须暂停或调度普通访问 \\
\bottomrule
\end{longtable}

\section{二、地址译码定义整机地图}

CPU 输出的是地址数值，整机必须规定这个数值由哪个目标响应。这个规定叫地址映射。一个简单教学板可以把低地址映射到 ROM，中间映射到 RAM，高地址小窗口映射到 GPIO、定时器和串口等设备寄存器。地址译码器把高位地址和访问类型变成片选信号。

\begin{lstlisting}
example 16-bit address map:
  0x0000-0x7FFF  boot ROM
  0x8000-0xEFFF  SRAM
  0xF000         GPIO_OUT
  0xF004         GPIO_IN
  0xF010         TIMER_COUNT
  0xF014         TIMER_CTRL
\end{lstlisting}

\begin{pdfdiagram}
  \node[hw state, minimum width=2.0cm] (addr) at (0,0) {地址 A[15:0]};
  \node[hw compute, minimum width=2.15cm] (dec) at (3.0,0) {高位译码};
  \node[hw state, minimum width=1.6cm] (rom) at (6.0,1.05) {ROM CS};
  \node[hw state, minimum width=1.6cm] (ram) at (6.0,0) {RAM CS};
  \node[hw control, minimum width=1.6cm] (io) at (6.0,-1.05) {IO CS};
  \draw[hw bus] (addr) -- (dec);
  \draw[hw bus] (dec.east) |- (rom.west) node[pos=0.3,above,hw label]{A15=0};
  \draw[hw bus] (dec) -- (ram) node[pos=0.55,above,hw label]{10xx};
  \draw[hw bus] (dec.east) |- (io.west) node[pos=0.3,below,hw label]{1111};
  \node[hw label, text width=7.8cm] at (3.4,-1.95) {验收目标：任意一个地址最多选中一个目标；未映射区域也要有错误、固定值或超时规则。};
\end{pdfdiagram}
\diagramnote{地址译码不是软件表格，而是片选电路。片选重叠会造成总线争用或误写设备。}

\begin{longtable}{L{0.18\linewidth}L{0.38\linewidth}L{0.34\linewidth}}
\toprule
映射对象 & 响应规则 & 典型错误 \\
\midrule
ROM & 复位入口和只读代码地址命中 & 复位后 CPU 取不到第一条指令 \\
RAM & 普通数据读写地址命中 & 写使能误打到 ROM 或设备窗口 \\
GPIO & 设备寄存器地址命中并产生副作用 & 地址译码过宽，普通内存写误触发外设 \\
未映射区 & 返回错误、超时或固定空值 & 访问悬空，CPU 永久等待或读到随机总线值 \\
\bottomrule
\end{longtable}

\topic{片选表达式就是硬件合同}
若地址空间为 64 KiB，可以用高位地址产生粗粒度片选。例如 \code{A15=0} 选 ROM，\code{A15=1,A14=0} 选 RAM，\code{A15=1,A14=1,A13=1,A12=1} 选 I/O 窗口。无论用门电路、译码器、CPLD 还是 FPGA 实现，都必须证明任意合法周期最多只有一个目标片选有效。

\begin{lstlisting}
rom_cs = !A15
ram_cs = A15 && !A14
io_cs  = A15 && A14 && A13 && A12
\end{lstlisting}

这三个表达式看起来像逻辑练习，实际上就是整机安全边界。若 ROM 和 RAM 可同时片选，读周期可能出现两个芯片同时驱动数据线；若 RAM 和 GPIO 片选重叠，普通 STORE 可能既改内存又改外设输出；若未映射地址没有超时或错误返回，CPU 可能等待一个永远不会响应的目标。

\topic{时序预算决定是否需要等待状态}
地址译码需要时间，存储器访问需要时间，数据到达 CPU 输入端后还要满足建立时间。若这些时间总和超过时钟周期，就必须降低频率、换更快器件，或插入等待状态。等待状态不是软件延时，而是硬件事务延长：地址、片选、读写方向和数据驱动关系必须在等待期间保持可解释。

\begin{lstlisting}
read_budget =
  cpu_address_output_delay
  + address_decode_delay
  + memory_access_time
  + input_setup_time
  + board_margin
\end{lstlisting}

若 \code{read_budget} 大于一个周期，控制器就不能在下一拍采样数据。慢 ROM、慢外设和长板级走线都可能需要 READY/WAIT 机制。正式规格还应说明最大等待长度和超时行为，避免坏设备把整机永久挂住。

\begin{longtable}{L{0.2\linewidth}L{0.34\linewidth}L{0.36\linewidth}}
\toprule
处理方式 & 适合场景 & 代价 \\
\midrule
降低时钟 & 全系统都慢，教学板容易观察 & 所有操作都变慢 \\
更快器件 & SRAM 或译码器成为瓶颈 & 成本、功耗或电平兼容性变化 \\
插入等待 & 只有部分目标较慢 & 控制器必须保持事务状态 \\
超时返回错误 & 目标可能不存在或失效 & 需要错误状态和恢复入口 \\
\bottomrule
\end{longtable}

\section{三、MMIO 把设备寄存器放进地址空间}

设备寄存器可以像内存地址一样被 CPU 读写，但它们不是普通 RAM。读状态寄存器可能清除某些位，写命令寄存器可能启动外设动作，写数据寄存器可能进入 FIFO，读数据寄存器可能弹出一个元素。这类地址称为内存映射 I/O，简称 MMIO。MMIO 的关键不是地址长得像内存，而是每个寄存器都有副作用和顺序规则。

\begin{longtable}{L{0.18\linewidth}L{0.36\linewidth}L{0.36\linewidth}}
\toprule
设备寄存器 & 读写语义 & 验收问题 \\
\midrule
\code{DIR} & 设置 GPIO 引脚方向 & 输出引脚和输入引脚不会混用 \\
\code{OUT} & 写输出锁存器 & 写入只改变目标输出位 \\
\code{IN} & 读取同步后的外部输入 & 按键抖动和亚稳态不能直接进入 CPU \\
\code{STAT} & 返回 ready、busy、error、irq 等状态 & 清除规则不会丢掉同时到来的新事件 \\
\code{DATA} & 数据收发寄存器或 FIFO 端口 & 空读、满写和溢出有明确定义 \\
\bottomrule
\end{longtable}

\topic{MMIO 写响应不等于设备已经完成}
CPU 写一个命令寄存器，通常只表示设备已经接收命令。设备真正完成可能要等若干周期、若干微秒，甚至更久。完成可以通过状态位、轮询、DMA completion 或中断通知。若把 MMIO 写返回当作“设备动作完成”，软件会在设备尚未写完缓冲区、尚未刷新状态或尚未清除错误时继续执行，形成难复现的硬件错误。

\begin{lstlisting}
device command sequence:
  CPU writes COMMAND register
  bus write response returns
  device starts internal work
  device sets DONE or ERROR
  interrupt or polling tells CPU to inspect result
\end{lstlisting}

\begin{pdfdiagram}
  \node[hw state, minimum width=1.55cm] (cpu) at (0,0) {CPU};
  \node[hw compute, minimum width=1.75cm] (bus) at (2.3,0) {总线响应};
  \node[hw control, minimum width=1.95cm] (cmd) at (4.85,0) {命令寄存器};
  \node[hw control, minimum width=1.95cm] (dev) at (7.45,0) {设备状态机};
  \node[hw state, minimum width=1.5cm] (done) at (9.75,0) {DONE};
  \draw[hw bus] (cpu) -- node[above=2pt,hw label,fill=white,inner sep=1pt]{写命令} (bus);
  \draw[hw bus] (bus) -- node[above=2pt,hw label,fill=white,inner sep=1pt]{写已接受} (cmd);
  \draw[hw bus] (cmd) -- node[above=2pt,hw label,fill=white,inner sep=1pt]{开始工作} (dev);
  \draw[hw return] (dev) -- node[above=2pt,hw label,fill=white,inner sep=1pt]{完成/错误} (done);
  \node[hw label, text width=9.0cm] at (4.7,-1.0) {总线写响应只证明命令到达设备接口；真正完成必须看状态位、中断或 DMA completion。};
\end{pdfdiagram}
\diagramnote{MMIO 写命令的两层完成边界。总线完成不等于设备动作完成。}

\topic{端口 I/O 与 MMIO 是两种地址合同}
8086 家族保留了独立 I/O 地址空间，\code{IN} 和 \code{OUT} 指令访问端口；许多 RISC 和单片机系统则把设备寄存器映射进普通地址空间，用 LOAD/STORE 访问。两者都能控制设备，但硬件合同不同。端口 I/O 需要 I/O 周期和端口地址译码；MMIO 需要内存地址译码和内存属性规则，例如不可缓存、不可合并或必须保持顺序。

\topic{外部输入必须先调理和同步}
按键、传感器、串口线和外部中断引脚都不按 CPU 时钟整齐变化。按键还会抖动，异步输入还可能让触发器进入亚稳态。教学板上的 GPIO 输入至少应经过电平保护、上拉/下拉、去抖策略和同步触发器，再进入设备寄存器。否则 CPU 读到的不是稳定硬件事实，而是未定义的边界行为。

\begin{warningbox}
不要把设备寄存器当作普通变量。普通 RAM 的读写目标是保存数据；MMIO 的读写目标是驱动硬件状态机。编译器、cache、写缓冲和乱序执行都可能改变普通内存访问的时机，但设备访问通常需要不可缓存属性、顺序约束和明确的完成检查。
\end{warningbox}

\topic{从 GPIO、定时器和 UART 看设备寄存器}
设备寄存器最好用具体布局理解。一个 GPIO 控制器可以有输出寄存器、输入寄存器、方向寄存器和中断状态寄存器；定时器可以有计数值、比较值、控制位和中断清除位；UART 可以有发送数据、接收数据、状态和波特率配置寄存器。CPU 访问这些寄存器时走的是地址事务，但设备内部看到的是控制状态机输入。

\begin{longtable}{L{0.18\linewidth}L{0.36\linewidth}L{0.36\linewidth}}
\toprule
设备 & 寄存器边界 & 常见验收问题 \\
\midrule
GPIO & \code{DIR/OUT/IN/IRQ_STAT} & 输入是否同步，输出是否只改目标位 \\
定时器 & \code{COUNT/CMP/CTRL/IRQ_CLEAR} & 到期事件是否锁存，清除是否丢新事件 \\
UART & \code{TX/RX/STAT/BAUD} & 空读、满写、错误位和中断顺序是否清楚 \\
\bottomrule
\end{longtable}

\section{四、中断和 DMA 把外部事件接回 CPU}

轮询可以让 CPU 一直读状态寄存器，直到设备 ready。但若设备很慢，CPU 会浪费大量周期。中断提供另一种路径：设备在状态变化后请求中断，中断控制器汇总、屏蔽和排序请求，CPU 在可接受边界保存当前执行上下文，跳到中断入口，服务程序读取设备状态并清除事件。

\begin{longtable}{L{0.18\linewidth}L{0.38\linewidth}L{0.34\linewidth}}
\toprule
对象 & 硬件职责 & 验收问题 \\
\midrule
设备 & 产生完成、错误或数据到达事件 & 状态位和中断请求是否同步一致 \\
中断控制器 & 汇总、屏蔽、优先级和投递 & 屏蔽位、优先级、EOI 顺序是否正确 \\
CPU 入口 & 保存边界状态并跳到入口 & PC、标志和特权边界是否可恢复 \\
服务程序 & 读状态、搬数据、清除事件 & 不丢新事件，不重复进入同一旧事件 \\
\bottomrule
\end{longtable}

\topic{中断清除顺序是硬件合同}
若服务程序先清设备状态，再清中断控制器，或先通知中断控制器再清设备，效果可能不同。某些设备要求先读状态后读数据，某些要求写 1 清位，某些要求读数据寄存器自动清除 ready。教材级规格必须写清楚事件锁存、状态读取、清除、EOI 和重新开放中断的顺序。否则系统在低速测试中正常，在高频事件下丢中断或重复中断。

DMA 解决的是另一类浪费。若磁盘、网卡、显示控制器或音频设备需要搬运大量数据，CPU 不应逐字从设备读出再逐字写入内存。DMA 让设备或 DMA 控制器成为总线发起者，按描述符或寄存器配置直接读写主存。CPU 负责准备缓冲区和描述符，设备负责搬运，完成后用状态位或中断通知 CPU。

\begin{lstlisting}
DMA receive sketch:
  CPU allocates buffer
  CPU writes descriptor address and length
  CPU ensures descriptor is visible to device
  device reads descriptor and writes buffer
  device posts completion
  CPU handles interrupt and inspects buffer
\end{lstlisting}

\topic{DMA 最容易错在可见性和所有权}
DMA 让 CPU 和设备共享主存，也引入 cache 一致性、顺序和缓冲区所有权问题。CPU 写好的描述符必须先对设备可见；设备写入的缓冲区必须在 CPU 读取前变得可见；CPU 不能在设备仍拥有缓冲区时重用它；设备不能越过分配的地址范围。现代系统还会用 IOMMU 限制 DMA 能访问哪些物理页，避免坏设备或错误驱动破坏任意内存。

\begin{pdfdiagram}
  \node[hw state, minimum width=1.85cm] (cpu) at (0,0) {CPU owns};
  \node[hw compute, minimum width=2.1cm] (desc) at (2.65,0) {描述符可见};
  \node[hw control, minimum width=1.95cm] (dev) at (5.35,0) {Device owns};
  \node[hw compute, minimum width=2.0cm] (comp) at (7.95,0) {completion};
  \node[hw state, minimum width=1.85cm] (cpu2) at (10.45,0) {CPU owns};
  \draw[hw bus] (cpu) -- node[above,hw label]{提交描述符} (desc);
  \draw[hw bus] (desc) -- node[above,hw label]{设备搬运} (dev);
  \draw[hw bus] (dev) -- node[above,hw label]{写状态/中断} (comp);
  \draw[hw return] (comp) -- node[above,hw label]{同步 cache 后读取} (cpu2);
  \node[hw label, text width=10.0cm] at (5.25,-1.0) {所有权不能重叠：CPU 交出缓冲区后不得改写；设备 completion 到达且数据可见后，CPU 才能重新使用。};
\end{pdfdiagram}
\diagramnote{DMA 缓冲区所有权。完成通知、cache 可见性和所有权回收必须按顺序发生。}

\begin{longtable}{L{0.18\linewidth}L{0.38\linewidth}L{0.34\linewidth}}
\toprule
边界 & 必须说明 & 失败后果 \\
\midrule
描述符可见 & CPU 写描述符后如何让设备看到 & 设备读到旧地址、旧长度或旧标志 \\
数据可见 & 设备写缓冲区后 CPU 如何看到新数据 & CPU 读到 cache 中的旧内容 \\
所有权 & 缓冲区何时归 CPU，何时归设备 & 重用正在 DMA 的缓冲区，数据被覆盖 \\
地址权限 & 设备能访问哪些物理页或 I/O 虚拟地址 & DMA 写坏无关内存或触发 IOMMU fault \\
完成顺序 & completion 到达是否代表所有数据已可见 & 收到中断后仍读到未完成数据 \\
\bottomrule
\end{longtable}

\section{五、从教学总线到现代互连}

早期微机的并行总线比较直观：地址线、数据线、读写控制、片选和等待信号都能在逻辑分析仪上看到。现代芯片内部和 PCIe 这类外部互连更复杂，可能使用分层 packet、请求队列、completion、错误报告和流控。但抽象问题没有改变：事务由发起者提出，经互连路由到目标，目标返回数据或完成状态，错误路径必须被记录和上报。

\begin{longtable}{L{0.18\linewidth}L{0.36\linewidth}L{0.36\linewidth}}
\toprule
互连对象 & 旧总线中的类比 & 现代形式 \\
\midrule
地址阶段 & 地址线和片选 & 请求 packet、BAR、路由 ID、地址译码 \\
数据阶段 & 数据线和字节使能 & payload、byte enable、burst 或 cache line \\
等待 & READY/WAIT & credit、ready/valid、队列背压 \\
完成 & 采样数据或写周期结束 & completion packet、状态位、中断或 MSI \\
错误 & 总线错误或超时 & poison、fault、AER、IOMMU fault、设备 error \\
\bottomrule
\end{longtable}

\topic{valid/ready 是理解现代片上互连的最小模型}
许多片上接口都可以抽象成 valid/ready 握手。发送方在 valid 为 1 时保持请求或响应稳定，接收方在 ready 为 1 时表示可以接收，只有二者同周期为 1，事务片段才真正完成。请求通道和响应通道可以分开，读请求和读数据也可以隔很多周期。这个模型比“总线一拍完成”更接近现代硬件。

\begin{pdfdiagram}[x=1cm,y=1cm]
  \draw[BookRule,line width=0.45pt] (0,1.2) -- (9.4,1.2);
  \draw[BookRule,line width=0.45pt] (0,0.45) -- (9.4,0.45);
  \draw[BookRule,line width=0.45pt] (0,-0.3) -- (9.4,-0.3);
  \node[left,hw label] at (0,1.2) {valid};
  \node[left,hw label] at (0,0.45) {ready};
  \node[left,hw label] at (0,-0.3) {payload};
  \draw[BookTeal,line width=0.85pt] (1.0,1.2) -- (7.6,1.2);
  \draw[BookAccent,line width=0.85pt] (4.0,0.45) -- (8.3,0.45);
  \draw[BookMuted,line width=0.65pt] (1.0,-0.3) -- (8.3,-0.3);
  \draw[BookRule,line width=0.45pt] (4.0,1.55) -- (4.0,-0.65);
  \node[hw label] at (4.0,1.8) {transfer};
  \node[hw label, text width=8.4cm] at (4.7,-1.05) {valid=1 且 ready=0 时，发送方必须保持 payload；二者同时为 1 的周期才完成一次传输。};
\end{pdfdiagram}
\diagramnote{valid/ready 握手。背压不是错误，而是接收方暂时不能接受。}

\begin{lstlisting}
handshake rule:
  transfer happens only when valid == 1 and ready == 1
  while valid == 1 and ready == 0:
    sender keeps payload stable
  receiver may apply backpressure by deasserting ready
\end{lstlisting}

\topic{写缓冲和内存顺序解释了现代芯片为什么看起来更快}
写入比读取更容易被隐藏。CPU 可以先把写事务放进写缓冲，让执行流水线继续推进；写缓冲再把相邻写合并、排队发送到 cache 或互连。这样常见 STORE 的表面延迟很低，但代价是可见性变复杂：其他核心、DMA 或设备什么时候看到这次写，取决于 cache 一致性、内存类型和屏障规则。普通 RAM、MMIO 和 DMA 描述符不能使用同一套可见性假设。

\topic{总线错误和超时必须成为规格的一部分}
未映射地址不能只写成“不要访问”。真实系统必须规定访问未映射区域时会发生什么：返回固定值、保持等待直到超时、产生总线错误，或进入异常入口。设备无响应、DRAM 控制器报错、ECC 检测失败、外设权限不允许，也都应该形成可观察错误状态。没有错误合同的整机，调试时会把硬件故障误判为程序死循环。

\topic{低时延来自近处、缓存、并行和卸载}
现代整机降低时延，不是让所有路径都变成零等待。L1/L2 cache 让常见数据走近处 SRAM；TLB 让近期地址转换不必每次查页表；预取和分支预测提前准备可能需要的路径；流水线、旁路和乱序调度让独立工作重叠；DMA、队列和中断聚合把批量 I/O 从 CPU 指令流中卸载。失败路径仍然存在：cache miss、TLB miss、DRAM 排队、设备背压、DMA fault 和错误中断都会把访问拉回慢路径。

\begin{classicnote}
读任何硬件平台，都可以用同一张事务表约束理解：发起者、目标、地址、方向、数据宽度、握手、完成、错误、副作用和可观察证据。早期 8086 总线、微控制器片内外设、PCIe 设备和现代 SoC 互连的规模不同，但这十项问题不变。
\end{classicnote}

\section*{章末练习与验收任务}

\begin{longtable}{L{0.22\linewidth}L{0.5\linewidth}L{0.18\linewidth}}
\toprule
任务 & 要求 & 验收标准 \\
\midrule
SRAM 读写周期 & 画出地址、CS、OE、WE、数据线和采样窗口 & 能指出谁驱动数据线，谁采样 \\
地址映射 & 设计 64 KiB ROM/RAM/GPIO/定时器映射 & 任意地址最多一个片选有效 \\
等待状态 & 为慢 ROM 和慢 I/O 设计 READY/WAIT 规则 & 等待期间地址和方向稳定，有超时策略 \\
MMIO 寄存器 & 设计 GPIO 的 \code{DIR/OUT/IN/STAT} 语义 & 外部输入经过同步，状态清除规则明确 \\
中断顺序 & 写出设备 ready 到 CPU 服务程序返回的 trace & 不丢事件，不重复清旧事件 \\
DMA 顺序 & 写出描述符、缓冲区、completion 和 cache 可见性边界 & CPU 和设备所有权不重叠 \\
valid/ready & 用 valid/ready 描述一次读请求和读响应 & 能说明背压和 completion 的作用 \\
未映射访问 & 规定未映射地址的超时、错误或固定返回值 & CPU 不会永久等待未知目标 \\
\bottomrule
\end{longtable}

\section*{参考解答与验收要点}

\begin{longtable}{L{0.18\linewidth}L{0.48\linewidth}L{0.24\linewidth}}
\toprule
任务 & 参考解答 & 必须保留的验收点 \\
\midrule
SRAM 周期 & 读时地址和片选稳定，OE 有效、WE 无效，由 SRAM 驱动数据；写时 OE 无效，由发起者驱动数据，WE 在地址和数据稳定后有效。 & 输出使能互斥，写窗口满足建立/保持时间。 \\
地址映射 & 例如低半区 ROM，高半区中一块 RAM，最高 4 KiB 为 I/O；用高位地址产生互斥片选，未映射区返回错误或超时。 & 片选互斥和未映射行为必须写入规格。 \\
等待状态 & 慢目标在 ready 之前保持事务未完成，控制器保持地址、片选和方向；超过最大等待则返回错误。 & 等待状态是硬件事务延长。 \\
MMIO & GPIO 输出写入锁存器，输入来自同步后的外部引脚，状态位说明 ready/error/irq，清除规则避免新旧事件混淆。 & MMIO 有副作用，不是普通变量。 \\
中断 & 设备锁存事件并请求中断，中断控制器投递，CPU 进入入口，服务程序读状态、处理数据、按设备规则清事件，再发送 EOI 或恢复中断。 & 清除顺序决定是否丢事件。 \\
DMA & CPU 准备描述符和缓冲区并保证对设备可见，设备搬运数据并投递 completion，CPU 在 completion 后按一致性规则读取缓冲区。 & 描述符可见、数据可见和缓冲区所有权必须分开。 \\
valid/ready & 只有 valid 和 ready 同时为 1 才传输；ready 为 0 时发送方保持 payload；读请求和读响应可隔多个周期。 & 背压不等于错误，completion 才是完成证据。 \\
未映射访问 & 未映射区域可以返回固定值、产生错误或等待到超时，但规则必须固定且可观察；不能让总线悬空成为随机值。 & 错误路径也是硬件规格。 \\
\bottomrule
\end{longtable}
